\chapter{Vision commerciale de MABRID}
\label{ch:business-vision}

\section{Introduction}

La place de marché numérique est devenue une institution déterminante de l'économie du vingt-et-unième siècle. Les plateformes qui mettent en relation acheteurs et vendeurs à grande échelle génèrent des effets de réseau, réduisent les coûts de recherche et démocratisent l'accès aux opportunités commerciales. Sur les marchés émergents, toutefois, la réussite d'une place de marché dépend autant de la \emph{confiance} que de la technologie : les consommateurs doivent croire que les annonces sont authentiques, que les litiges seront traités et que les données personnelles seront gérées de manière responsable. Ce chapitre établit la justification commerciale de MABRID, une place de marché numérique marocaine de confiance, en articulant la vision, la mission, la valeur pour les parties prenantes et des objectifs commerciaux mesurables, sans référence aux mécanismes d'implémentation.

\section{Économie des places de marché numériques}

\subsection{Économie des plateformes et marchés multi-acteurs}

Une place de marché numérique est une plateforme multi-acteurs qui crée de la valeur en facilitant les interactions entre des groupes d'utilisateurs distincts---généralement acheteurs et vendeurs---dont la participation est interdépendante \parencite{rochet2003platform,evans2016matchmakers}. Contrairement aux entreprises linéaires traditionnelles, les opérateurs de plateformes investissent dans les règles de gouvernance, les systèmes de réputation et les mécanismes de découverte qui accroissent l'efficacité des transactions dans l'écosystème.

Les économistes distinguent les effets de réseau \emph{intra-groupe} (davantage d'acheteurs attirent davantage d'acheteurs grâce aux avis partagés) et les effets \emph{inter-groupes} (davantage de vendeurs attirent davantage d'acheteurs et inversement). Le modèle économique de MABRID suppose des effets inter-groupes positifs modérés par des contrôles de qualité : une croissance illimitée sans vérification éroderait la confiance et réduirait finalement la participation.

\figureplaceholder{H}{figures/ch01-platform-ecosystem.pdf}{0.85}{Écosystème d'une place de marché multi-acteurs : acheteurs, vendeurs, gouvernance de la plateforme et parties prenantes externes}

\subsection{La confiance comme actif stratégique}

Les déficits de confiance demeurent un obstacle majeur à l'adoption du commerce numérique dans de nombreuses régions. La recherche sur la confiance dans les plateformes met en évidence le rôle des signaux institutionnels---badges de vérification, politiques transparentes, modération réactive---dans la réduction du risque perçu \parencite{mcknight2002trust,eisenmann2006strategies}. MABRID positionne la confiance non comme un slogan marketing, mais comme une capacité opérationnelle soutenue par la gouvernance, la sécurité et la conformité juridique (développées dans les chapitres suivants).

\subsection{Étapes du cycle de vie d'une place de marché}

Les places de marché évoluent généralement à travers le lancement (acquisition de l'offre), la croissance (liquidité et engagement), la maturité (monétisation et différenciation) et l'expansion (mise à l'échelle géographique ou verticale) \parencite{chen2018liquidity}. La stratégie de lancement de MABRID privilégie la \emph{liquidité} dans les villes d'ancrage avant le déploiement national, afin que les acheteurs trouvent une offre de qualité suffisante et que les vendeurs bénéficient d'une visibilité significative.

\section{La transformation numérique marocaine}

\subsection{Contexte macroéconomique et politique}

Le Maroc a poursuivi sa transformation numérique par des stratégies nationales reliant l'investissement dans les infrastructures, le soutien à l'entrepreneuriat et la modernisation administrative \parencite{mcit2020digital,morocco2025digital}. La pénétration mobile, l'amélioration de la couverture haut débit et la familiarité croissante avec les paiements numériques créent des conditions favorables aux plateformes de place de marché---à condition qu'elles respectent les attentes réglementaires locales et les normes culturelles de construction des relations commerciales.

\subsection{Écart de numérisation des PME}

Les petites et moyennes entreprises (PME) représentent une part substantielle de l'emploi au Maroc, mais disposent souvent d'une présence numérique professionnelle limitée au-delà de pages informelles sur les réseaux sociaux. MABRID répond à l'\emph{écart de découvrabilité} : aider les entreprises légitimes à devenir trouvables, évaluables et joignables sans exiger de chaque commerçant qu'il construise une infrastructure numérique indépendante.

\subsection{Urbanisation et communautés par ville}

Le commerce marocain reste fortement orienté vers les villes. Une plateforme nationale doit refléter la pertinence géographique---les utilisateurs à Casablanca, Rabat, Marrakech et dans les villes secondaires attendent des résultats localement significatifs. Le modèle de communautés par ville de MABRID (gouverné par des politiques de plateforme plutôt que par des détails techniques) favorise l'engagement hyperlocal tout en préservant la cohérence de la marque nationale.

\figureplaceholder{H}{figures/ch01-morocco-map.pdf}{0.75}{Couverture nationale prévue : villes d'ancrage et corridors d'expansion par phases}

\section{Opportunité commerciale}

\subsection{Énoncé du problème}

Les acheteurs peinent à découvrir des entreprises locales fiables à travers des canaux fragmentés. Les vendeurs font face à des coûts d'acquisition client élevés et à des outils limités pour démontrer leur crédibilité. Aucun des deux groupes ne bénéficie d'intermédiaires opaques ou d'environnements d'annonces non modérés. MABRID intervient comme couche de découverte sélectionnée et de confiance.

\subsection{Logique de dimensionnement du marché}

L'opportunité de marché peut s'exprimer par le \gls{tam}, le \gls{sam} et le \gls{som} \parencite{osterwalder2010business} :
\begin{itemize}
  \item \textbf{TAM :} Dépenses totales des consommateurs et des PME adressables par la découverte locale et la génération de prospects à l'échelle nationale.
  \item \textbf{SAM :} Populations urbaines et connectées utilisant activement les smartphones pour la recherche commerciale.
  \item \textbf{SOM :} Part réaliste capturable dans les premières années compte tenu de la dynamique concurrentielle et de la capacité de mise sur le marché.
\end{itemize}

Le tableau~\ref{tab:market-sizing} présente des hypothèses de dimensionnement illustratives pour la planification stratégique. Les chiffres doivent être validés périodiquement à l'aide des statistiques nationales et de recherches primaires.

\begin{table}[H]
  \centering
  \caption{Cadre illustratif de dimensionnement du marché pour MABRID (hypothèses de planification)}
  \label{tab:market-sizing}
  \begin{tabularx}{\textwidth}{@{}lXr@{}}
    \toprule
    \textbf{Segment} & \textbf{Description} & \textbf{Échelle indicative} \\
    \midrule
  TAM & Découverte du commerce local national et présence numérique des PME & Élevée \\
  SAM & Acheteurs urbains équipés de smartphones et PME enregistrées dans les villes de lancement & Moyenne \\
  SOM & Utilisateurs actifs et vendeurs payants de MABRID sous 36 mois & Ciblée \\
    \bottomrule
  \end{tabularx}
\end{table}

\subsection{Paysage concurrentiel}

Les concurrents comprennent les plateformes généralistes de petites annonces, le commerce social, les annuaires sectoriels et les réseaux informels de commerce via WhatsApp. MABRID se différencie par des profils d'entreprise vérifiés, des avis structurés, des communautés par ville, des programmes vendeurs premium et des opérations de confiance et sécurité de niveau entreprise---des domaines où les canaux informels ne peuvent reproduire une crédibilité institutionnelle.

\section{Vision et mission}

\subsection{Vision}

\textbf{Devenir la place de marché numérique la plus fiable du Maroc pour découvrir et interagir avec les entreprises locales.}

\subsection{Mission}

\textbf{Donner aux acheteurs et aux vendeurs les moyens d'agir grâce à une découverte transparente, des signaux de confiance vérifiables et une gouvernance responsable de la plateforme respectant le droit marocain et les meilleures pratiques internationales en matière de protection des données et de sécurité.}

\subsection{Piliers stratégiques}

\begin{enumerate}
  \item \textbf{Confiance par conception} --- vérification, modération et politiques claires.
  \item \textbf{Pertinence locale} --- découverte centrée sur la ville et engagement communautaire.
  \item \textbf{Participation inclusive} --- intégration accessible pour les micro-entreprises et les artisans.
  \item \textbf{Excellence opérationnelle} --- sécurité, conformité et continuité des activités comme fondations non négociables.
\end{enumerate}

\section{Proposition de valeur}

\subsection{Valeur pour l'acheteur}

Les acheteurs bénéficient d'une recherche consolidée par catégories (alimentation, services, santé, automobile, services professionnels), d'une découverte cartographique, d'avis authentiques, d'options de contact direct et de recommandations communautaires au sein de leur ville. Les abonnements premium acheteurs peuvent offrir une personnalisation renforcée et un support prioritaire sans modifier l'accès gratuit de base.

\subsection{Valeur pour le vendeur}

Les vendeurs disposent d'une vitrine professionnelle, de la gestion des annonces, de canaux de demande, d'analyses sur la visibilité et l'engagement, d'un placement en vedette optionnel et de parcours de vérification signalant la légitimité. La plateforme réduit la dépendance à la seule publicité sociale payante.

\subsection{Valeur pour la plateforme}

MABRID capture de la valeur par des niveaux d'abonnement, des services de vérification, des partenariats d'entreprise et des revenus auxiliaires futurs (sous réserve de la permissibilité réglementaire). La plateforme n'intermédie pas initialement les paiements entre acheteurs et vendeurs, ce qui réduit la complexité réglementaire financière tout en se concentrant sur la génération de prospects et la réputation.

\section{Parties prenantes cibles}

\subsection{Acheteurs}

Consommateurs individuels et ménages recherchant des entreprises locales fiables. Les segments comprennent les professionnels urbains, les familles, les expatriés et les touristes nationaux. La confiance, la commodité et une expérience mobile d'abord favorisent l'adoption.

\subsection{Vendeurs}

Micro-entreprises, PME, franchises et professionnels indépendants. L'intégration doit minimiser les frictions tout en préservant des options de vérification pour les catégories à haute confiance (santé, juridique, conseil financier).

\subsection{Entreprises et partenaires}

Les banques, opérateurs télécoms, associations professionnelles et chambres de commerce peuvent s'associer pour le co-marketing, la vérification ou des services numériques groupés.

\subsection{Gouvernement et régulateurs}

Les autorités publiques bénéficient d'une visibilité formalisée des PME, d'un alignement sur la protection des consommateurs et d'une activité économique imposable canalisée par des entreprises documentées. L'engagement auprès de la \gls{cndp} et des régulateurs sectoriels est essentiel.

\subsection{Investisseurs}

Les investisseurs en capital-risque et à impact évaluent la croissance du TAM, l'économie unitaire, la défendabilité (avantage concurrentiel fondé sur la confiance) et le risque réglementaire. Ce rapport soutient la diligence raisonnable sur la gouvernance et la préparation à la conformité.

\figureplaceholder{H}{figures/ch01-stakeholder-map.pdf}{0.9}{Carte des parties prenantes et principaux échanges de valeur}

\section{Objectifs commerciaux}

\subsection{Objectifs à court terme (0--12 mois)}

\begin{itemize}
  \item Atteindre la préparation au lancement dans les villes d'ancrage avec une cohorte de vendeurs vérifiés et une base d'acheteurs active.
  \item Établir les opérations de confiance et sécurité et publier la suite de politiques juridiques.
  \item Atteindre une base de sécurité alignée sur les fonctions Identifier--Protéger du \gls{nistcsf}.
  \item Sécuriser les premiers revenus issus des abonnements vendeurs et des services de vérification.
\end{itemize}

\subsection{Objectifs à moyen terme (12--36 mois)}

\begin{itemize}
  \item S'étendre à vingt villes marocaines ou plus avec des modérateurs communautaires locaux.
  \item Atteindre une marge de contribution positive sur les coûts variables cloud et support.
  \item Obtenir une évaluation de sécurité externe et progresser vers une feuille de route de certification \gls{iso27001}.
  \item Lancer un programme de partenariat d'entreprise pour les entreprises multi-sites.
\end{itemize}

\subsection{Objectifs à long terme (36+ mois)}

\begin{itemize}
  \item Devenir la marque de référence pour la découverte locale au Maroc.
  \item Évaluer une expansion sélective au Maghreb sous réserve de faisabilité juridique et opérationnelle.
  \item Introduire des services de confiance avancés (assurance d'identité, partenariats de médiation des litiges).
\end{itemize}

\subsection{Indicateurs de succès}

Le tableau~\ref{tab:business-kpis} relie les objectifs à des \gls{kpi} illustratifs. Les définitions opérationnelles et les cibles relèvent des documents internes de gestion de la performance.

\begin{table}[H]
  \centering
  \caption{Indicateurs commerciaux illustratifs pour MABRID}
  \label{tab:business-kpis}
  \begin{tabularx}{\textwidth}{@{}lXl@{}}
    \toprule
    \textbf{Dimension} & \textbf{KPI} & \textbf{Horizon} \\
    \midrule
    Croissance & Acheteurs et vendeurs actifs mensuels & Mensuel \\
    Confiance & Part de vendeurs vérifiés ; délai de résolution des litiges & Trimestriel \\
    Engagement & Annonces consultées par session ; taux de réponse aux messages & Hebdomadaire \\
    Revenus & ARPU ; conversion d'abonnement ; attrition & Mensuel \\
    Conformité & Achèvement de la formation aux politiques ; clôture des constats d'audit & Trimestriel \\
    \bottomrule
  \end{tabularx}
\end{table}

\section{Conclusion du chapitre}

La vision commerciale de MABRID répond à un besoin de marché documenté : une découverte locale de confiance dans une économie numérique marocaine en transformation. En articulant une valeur claire pour les parties prenantes, des piliers stratégiques et des objectifs mesurables, la plateforme établit une base pour la conception organisationnelle, l'investissement en sécurité et l'engagement réglementaire développés dans les chapitres suivants. L'accent mis sur la confiance et la gouvernance distingue MABRID des modèles de petites annonces axés sur le volume et aligne le projet sur une économie de place de marché durable et digne d'investissement.
